\section{Introduction}
Category theory was first introduced by MacLane and Eilenberg to give a precise definition to what it meant for something to be ``natural'' \cite{riehl}. The definition of naturality requires natural transformations, which in turn require functors, which reqiure categories. And thus, category theory was born. In this sense, category theory itself can be thought of as ``natural''.

Category theory has evolved over the near century since it's inception, and nowadays is of fundamental importance to many branches of mathematics, especially algberaic ones such as algebraic topology, algebraic geometry, and of course, algebra itself. This makes motivating it with algebraic and topological definitions very easy, and I will make use of these fields at an undergraduate level as motivation and examples throughout these notes. However, I don't want to assume extensive prerequisite knowledge in these fields, and I would prefer these notes to be accessible to anyone with only basic knowledge of what sets and functions are. Whether it's prudent or not to introduce students to categories before algebra is a debate\footnote{It's not.} for another day. My goal is not to teach a course on this material, but rather to give a friendly introduction, and as such I believe it fit to keep these notes as accessible as possible.

These notes will be structured into a few sections. I'll start with motivation, then will define a category and give some examples. After that I will continue with more important categorical definitions, along with a brief intermission for the Yoneda lemma. I'll finish up with some more advanced topics that I either think are important, or am personally interested in.
